\documentclass[10pt,compress]{beamer}
%\documentclass[fleqn,xcolor=dvipsnames]{beamer}
\usetheme{Madrid}
%\usetheme{Boadilla}
%\usetheme{default}
%\usetheme{Warsaw}
%\usetheme{CambridgeUS}
%\usetheme{Sybila}
%\usetheme[hideothersubsections]{Berkeley}
%\usetheme[hideothersubsections]{PaloAlto}
%%\usetheme[hideothersubsections]{Goettingen}
%\usetheme{CambridgeUS}
%\usetheme{Bergen} % This template has nagivation on the left
%\usetheme{Frankfurt} % Similar to the default with an extra region at the top.

% Uncomment the following line if you want page numbers and using Warsaw theme%
%\setbeamertemplate{footline}[page number]

\definecolor{aggiemaroon}{RGB}{80,0,0} % Official RGB code for aggie maroon
\definecolor{darkergray}{RGB}{28,28,28}
\definecolor{EconomistBlue}{RGB}{0, 38, 80}
\definecolor{EconomistBlueoth}{RGB}{2, 38, 77}
\usecolortheme[named=EconomistBlue]{structure}
\useoutertheme{shadow}
\useinnertheme{rounded}
\setbeamertemplate{navigation symbols}{}
\setbeamerfont{structure}{family=\rmfamily,series=\bfseries}
\usefonttheme[stillsansseriftext]{serif}

\usepackage{amsmath}
\usepackage{caption}
\usepackage{graphicx,comment}
\usepackage[english]{babel}
\usepackage{graphicx}
\usepackage{rotating}
\usepackage{multicol}
\usepackage{hyperref}
\usepackage{enumerate}
\usepackage{tikz}
\usepackage{graphicx}
\usepackage{bm}
\usepackage{subcaption}
\useoutertheme{miniframes}
\setbeamercolor{section in head/foot}{fg=white,bg=EconomistBlueoth}
\setbeamerfont*{section in head/foot}{family=\rmfamily,size=\fontsize{6}{8}}

%\setbeamercolor{subsection in head/foot}{fg=white,bg=black}


%\usebackgroundtemplate{
%\tikz\node[opacity=0.025, rotate = 0] {\includegraphics[height=\paperheight,width=\paperwidth]{TAM-Logo.png}};}
%{\includegraphics[height=1in,width=1in]{TAM-Logo.png}};}



\title[Texas A\&M University]{Anchor Institutions as Macroeconomic Stabilizers:\\The Case of College Towns}

\author [Sanin-Restrepo]{Sebastian Sanin-Restrepo}

\date[ \today]{\today}

\institute[Texas A\&M U.] % (optional, but mostly needed)
{Public Economics \\
 %\includegraphics[scale=0.60]{TAM-Logo1.png} \\ [0.0 cm]
 }




\begin{document}

\begin{frame}
\titlepage
\end{frame}

%\begin{frame}
%\frametitle{Overview} 
%\tableofcontents
%\end{frame}

%-----------------------------------------------------
\section{Motivation}
%-----------------------------------------------------

\begin{frame}{Why Study College Towns?}
	\begin{itemize}
		\item Universities are large, immobile “anchor institutions” with stable payrolls and countercyclical enrollment.
		\item They generate strong local linkages: non-tradables, housing, services.
		\item Yet dependence creates a \textbf{stability-fragility tradeoff}: resilience in normal recessions, vulnerability when campus presence collapses.
		\item COVID-19 exposed fragility: remote instruction shut down campus economies.
		\item Goal: \textbf{measure how university dependence shapes local cyclical sensitivity and asymmetric responses across recessions}.
	\end{itemize}
\end{frame}


\begin{frame}{Conceptual Background}
	\begin{itemize}
		\item Literature documents:
		\begin{itemize}
			\item Long-run development effects of universities (Abel \& Deitz 2012; Berlingieri et al. 2022; Ferhat 2022).
			\item Countercyclical enrollment + stable payrolls as stabilizers (Betts \& McFarland 1995; Dellas \& Koubi 2003).
			\item Universities as local demand anchors (Harker et al. 2024; Howard et al. 2024).
		\end{itemize}
		\item Closest work: Jump \& Scavette (2024) show flagships stabilize in 2008, but not in 2020.
		\item This project extends these results:
		\begin{itemize}
			\item Beyond flagship universities,
			\item Multiple recessions,
			\item Labor, housing, and non-tradables,
			\item Continuous exposure measure.
		\end{itemize}
		\item Relevance:
			\begin{itemize}
			\item Policy: Many regions are doubling down on anchor-based development, cities want universities or hospitals.
			\item Local public finance / planning: If reliance stabilizes usual recessions but amplifies rare shutdowns, towns should plan for that asymmetry.
			\end{itemize}
	\end{itemize}
\end{frame}

%-----------------------------------------------------
\section{Research Questions}
%-----------------------------------------------------

\begin{frame}{Research Questions}
	\begin{enumerate}
		\item Do university-dependent cities experience smaller declines in employment, payroll, and housing conditions during recessions?
		\item Does greater university exposure attenuate or amplify cyclical sensitivity?
		\item How did the 2020 shutdown alter the stabilizing mechanism?
		\item What roles do housing tightness, service-sector exposure, and commuting networks play?
	\end{enumerate}
\end{frame}

\begin{frame}{Key Hypotheses}
	\begin{itemize}
		\item \textbf{H1 (Stabilization):} In ordinary recessions, high university dependence reduces employment losses and unemployment spikes.
		\item \textbf{H2 (Fragility):} When in-person campus activity collapses, the stabilizing channel reverses.
		\item \textbf{H3 (Heterogeneity):} Stabilization is stronger in places with:
		\begin{itemize}
			\item Larger service sectors,
			\item Tighter housing markets,
			\item Higher commuting inflows connected to the university.
		\end{itemize}
	\end{itemize}
\end{frame}

%-----------------------------------------------------
\section{Data}
%-----------------------------------------------------

\begin{frame}{Data Overview}
	\textbf{Time frame:} 1990–2023 (subject to availability). \\
	\vspace{0.2cm}
	\begin{itemize}
		\item \textbf{University exposure:} IPEDS (NCES) - FTE enrollment, payroll, employment (NAICS 61/6113).
		\item \textbf{Local labor markets:} BLS QCEW - total and sectoral employment/payroll.
		\item \textbf{Unemployment:} BLS LAUS (county and selected cities).
		\item \textbf{Housing:} Zillow ZORI rents, FHFA house prices.
		\item \textbf{Demographics:} ACS + Decennial Census.
		\item \textbf{Commuting patterns:} LEHD LODES - worker residence/ workplace flows.
		\item \textbf{Recession periods:} NBER peak–trough dates.
	\end{itemize}
\end{frame}

\begin{frame}{Measuring University Dependence}
	\begin{itemize}
		\item \textbf{Continuous exposure measure}:
		\[
		s_i = \frac{\text{Higher-ed payroll or employment}}{\text{Total local payroll or employment}}
		\]
		\item \textbf{Alternative:} FTE students per capita.
		\item \textbf{Binary indicator:} “College town” if higher education $\geq$ 20–25\% of local employment.
		\item \textbf{Gumprecht (2008)}:``\textit{I consider as a college town any city where a college or university and the cultures it creates exert a dominant influence over the character of the town."}. Further measure $\approx$ 20\% pop. are 4-year college students and less than 350,000 people in the "city". (\textbf{This presentation.})
		\item Robust to alternative thresholds and continuous specifications.
	\end{itemize}
\end{frame}

\begin{frame}{Measuring University Dependence}
	\begin{figure}
		\centering
		\includegraphics[width=0.8\linewidth]{map1}
		\caption*{Figure 1: College Towns in US (Gumprecht, 2008, pg.5)}
		\label{fig:map1}
	\end{figure}
	
\end{frame}

%-----------------------------------------------------
\section{Empirical Strategy}
%-----------------------------------------------------

\begin{frame}{Identification Strategy}
	\begin{itemize}
		\item Recessions are national shocks; exposure to universities varies cross-sectionally.
		\item Use three identification approaches:
		\begin{enumerate}
			\item \textbf{Synthetic Difference-in-Differences (SDID)}
			\item \textbf{Interaction DiD + Local Projections}
			\item \textbf{Event study of 2020 shutdown}
		\end{enumerate}
		\item Core identification assumption:
		\begin{quote}
			After reweighting (SDID) or conditioning on fixed effects and pre-trends (DiD/LP), high- and low-exposure counties would follow similar paths absent recession.
		\end{quote}
	\end{itemize}
\end{frame}

\begin{frame}{SDID Framework}
	\begin{itemize}
		\item For each recession:
		\[
		Y_{it} = \alpha_i + \tau_t + \delta \cdot (D_i \times Recession_t) + \varepsilon_{it}
		\]
		\item SDID constructs unit and time weights $(\omega_i, \lambda_t)$ to exactly match pre-recession trends of treated and control counties.  
		\item Recovers causal effect of recessions on high-exposure counties.
		\item Implemented for 2001, 2008–09, and 2020.
	\end{itemize}
\end{frame}

\begin{frame}{Continuous Exposure DiD}
	\[
	Y_{it} = \alpha_i + \tau_t + \beta (Recession_t \times s_i) + X_{it}'\gamma + \varepsilon_{it}
	\]
	\begin{itemize}
		\item $\beta$ = marginal effect of exposure on recession sensitivity.
		\item Leads/lags used to test pre-trends.
		\item Add mechanisms:
		\[
		Y_{it} = \alpha_i + \tau_t + \beta_1 (Recession_t \times s_i) + \beta_2 (Recession_t \times s_i \times Z_i) + \varepsilon_{it}
		\]
		where $Z_i$ = housing elasticity, service sector share, etc.
	\end{itemize}
\end{frame}

\begin{frame}{Local Projections: Dynamic Adjustment}
	\begin{equation*}
		\Delta^h Y_{i,t_0+h} = \alpha_i + \kappa_h + \sum_{j=-L}^{H}\beta_{h,j}\,\mathbf{1}\{t=j\}\cdot s_i + X_{it_0}'\gamma_h + u_{i,t_0+h},
	\end{equation*}
	
	\begin{itemize}
		\item Horizon-specific dynamic responses across exposure levels.
		\item Compare depth and speed of recovery between high and low exposure counties.
	\end{itemize}
\end{frame}

\begin{frame}{Event Study: 2020 Shutdown}
	\[
	Y_{it} = \alpha_i + \tau_t + \sum_{k \neq -1} \delta_k (1\{t=k\} \cdot s_i) + X_{it}'\gamma + \varepsilon_{it}
	\]
	\begin{itemize}
		\item Captures the collapse of stabilizing channels when campus activity halted.
		\item Expect positive unemployment effects and negative payroll effects for high-exposure areas.
	\end{itemize}
\end{frame}

%-----------------------------------------------------
\section{Preliminary Evidence}
%-----------------------------------------------------

\begin{frame}{Unemployment Patterns (1990–2024)}
	\begin{itemize}
		\item College-town areas have lower median unemployment than non-college areas (county and MSA level).
		\item Gap narrows during recessions, suggesting modest stabilization.
		\item 2020 recession: effect reverses, consistent with fragility.
	\end{itemize}
	\begin{figure}[H]
		\centering
		\begin{subfigure}{0.48\textwidth}
			\centering
			\includegraphics[width=\linewidth]{../output/g1county}
			\caption{County level}
			\label{fig:g1_county}
		\end{subfigure}
		\begin{subfigure}{0.48\textwidth}
			\centering
			\includegraphics[width=\linewidth]{../output/g1msa}
			\caption{MSA level}
			\label{fig:g1_msa}
		\end{subfigure}
		\caption*{\textbf{Figure 1}: Unemployment by Area and College Town classification}
		\label{fig:g1_unemployment}
	\end{figure}
	
\end{frame}

\begin{frame}{Event Study Results}
	\begin{itemize}
		\item County-level: larger unemployment spikes in college towns in full sample, but effect disappears excluding 2020.
	\end{itemize}
	\begin{figure}[H]
		\centering
		\begin{subfigure}{0.48\textwidth}
			\centering
			\includegraphics[width=\linewidth]{../output/g2county_all}
			\caption{Counties – full sample}
			\label{fig:g2county_all}
		\end{subfigure}
		\hfill
		\begin{subfigure}{0.48\textwidth}
			\centering
			\includegraphics[width=\linewidth]{../output/g2county_pre}
			\caption{Counties – Pre-Covid sample}
			\label{fig:g2county_pre}
		\end{subfigure}
		\caption*{\textbf{Figure 2}: Event Study: Effect of Recessions on Unemployment in College Towns}
		\label{fig:g2_event_county}
	\end{figure}
	
\end{frame}

\begin{frame}{Event Study Results}
	\begin{itemize}
		\item MSA-level (preferred): no difference in full sample; pre-COVID, college towns slightly more resilient.
		\item Interpretation: stabilizing effect exists but weak—masked or reversed by 2020.
	\end{itemize}
	\begin{figure}[H]
		\centering
		\begin{subfigure}{0.48\textwidth}
			\centering
			\includegraphics[width=\linewidth]{../output/g2msa_all}
			\caption{MSAs – full sample}
			\label{fig:g2msa_all}
		\end{subfigure}
		\hfill
		\begin{subfigure}{0.48\textwidth}
			\centering
			\includegraphics[width=\linewidth]{../output/g2msa_pre}
			\caption{MSAs – Pre-Covid sample}
			\label{fig:g2msa_pre}
		\end{subfigure}
		\caption*{\textbf{Figure 2}: Event Study: Effect of Recessions on Unemployment in College Towns}
		\label{fig:g3_event_msa}
	\end{figure}
	
\end{frame}




%-----------------------------------------------------
\section{Next Steps}
%-----------------------------------------------------

\begin{frame}{Next Steps}
	\begin{itemize}
		\item Build full exposure measure using QCEW NAICS 6113 extraction.
		\item Complete SDID estimation for each recession.
		\item Expand to sectoral and housing outcomes.
		\item Incorporate commuting-based spatial spillovers (LODES).
		\item Draft theory section on stability-fragility tradeoff.
	\end{itemize}
\end{frame}











%\begin{frame}
%\frametitle{My model}
%\begin{flushleft}
%\begin{enumerate}
%\item Let $X_{i1},\ldots,X_{in}$, $i=1,\ldots,q$, are observations generated from the following model:
%\item  $X_{ij} \sim N(\mu_i, \sigma_i^2)$, $i=1,\ldots,q$, $j=1,\ldots,n$.
%\end{enumerate}
%\end{flushleft} 
%\end{frame}



%------------------------------------------------

\begin{frame}{}
  \centering \Huge
  \emph{Thank You!}
\end{frame}

%----------------------------------------------------------------------------------------
\end{document}



